% GENERATED FILE. Edit canonical lessons, then run npm run generate:books.
% canonical-source-hash: fnv1a64:31694517e0410c2a
% canonical-lessons: RU-C521-kist, RU-C521-kholst, RU-C521-akvarel, RU-C521-peyzazh, RU-C521-pozirovat

\chapter{Brush and Canvas}
\label{ch:ru-brush-and-canvas}
\hlchaptermodality{521}

\begin{chapteropening}
\textbf{By the end of this chapter:} \emph{I can talk about painting and sitting for a portrait.}

Complete the last lesson of chapter 521: I can talk about painting and sitting for a portrait.
\end{chapteropening}

\section[\texorpdfstring{\ru{кисть} (kist)}{кисть (kist)}]{\ru{кисть} (kist) — a paintbrush; the hand}

\label{lesson:RU-C521-kist}



\begin{quote}
Before the new one: say the Russian for a viewer, then the Russian for to act in a film.
\end{quote}

\subsection*{You'll want to know: \ru{кисть}}
\textbf{\ru{кисть}} (\emph{kist}) — \textquotedblleft{}a paintbrush; the hand\textquotedblright{}.

\begin{quote}
\ru{Художник} \ru{взял} \ru{тонкую} \ru{кисть}.
\end{quote}

\begin{quote}
\emph{Khudózhnik vzyal tónkuyu kist.}
\end{quote}

\begin{quote}
The artist picked up a fine brush.
\end{quote}

Feminine. Also the hand from the wrist down, and a bunch of grapes.

\begin{cousinweb}
EtymDB lists \emph{\ru{кисть}}, glossed 'cluster, bunch', without a source, and the origin is not settled.
\end{cousinweb}

\begin{grammarlens}[title={What you've built}]
One more word for telling, explaining and describing at B1.
\end{grammarlens}

\subsection*{Your turn}
\begin{itemize}
\raggedright
  \item \emph{Say it:} \emph{kist}
  \item \emph{Say it:} \emph{kist}, once more
  \item \emph{Say it:} a sentence of your own with \emph{kist}
  \item \emph{Recall:} say the Russian for a viewer, then the Russian for to act in a film, then say \emph{kist} again
\end{itemize}

\begin{tcolorbox}[breakable,colback=teal!4,colframe=teal!35!black,title={Before you move on}]
What does \ru{кисть} mean? (\textquotedblleft{}A paintbrush; the hand\textquotedblright{}.) Say it once more.
\end{tcolorbox}
\section[\texorpdfstring{\ru{холст} (kholst)}{холст (kholst)}]{\ru{холст} (kholst) — canvas}

\label{lesson:RU-C521-kholst}



\begin{quote}
Before the new one: say the Russian for to act in a film, then the Russian for a paintbrush.
\end{quote}

\subsection*{You'll want to know: \ru{холст}}
\textbf{\ru{холст}} (\emph{kholst}) — \textquotedblleft{}canvas\textquotedblright{}.

\begin{quote}
\ru{Картина} \ru{написана} \ru{маслом} \ru{на} \ru{холсте}.
\end{quote}

\begin{quote}
\emph{Kartína napísana máslom na kholsté.}
\end{quote}

\begin{quote}
The picture is painted in oils on canvas.
\end{quote}

Masculine; genitive \ru{холста}.

\begin{cousinweb}
EtymDB lists \emph{\ru{холст}} without a source, and the origin is not settled.
\end{cousinweb}

\begin{grammarlens}[title={What you've built}]
One more word for telling, explaining and describing at B1.
\end{grammarlens}

\subsection*{Your turn}
\begin{itemize}
\raggedright
  \item \emph{Say it:} \emph{kholst}
  \item \emph{Say it:} \emph{kholst}, once more
  \item \emph{Say it:} a sentence of your own with \emph{kholst}
  \item \emph{Recall:} say the Russian for to act in a film, then the Russian for a paintbrush, then say \emph{kholst} again
\end{itemize}

\begin{tcolorbox}[breakable,colback=teal!4,colframe=teal!35!black,title={Before you move on}]
What does \ru{холст} mean? (\textquotedblleft{}Canvas\textquotedblright{}.) Say it once more.
\end{tcolorbox}
\section[\texorpdfstring{\ru{акварель} (akvarél)}{акварель (akvarél)}]{\ru{акварель} (akvarél) — watercolour}

\label{lesson:RU-C521-akvarel}



\begin{quote}
Before the new one: say the Russian for a paintbrush, then the Russian for canvas.
\end{quote}

\subsection*{You'll want to know: \ru{акварель}}
\textbf{\ru{акварель}} (\emph{akvarél}) — \textquotedblleft{}watercolour\textquotedblright{}.

\begin{quote}
\ru{Она} \ru{рисует} \ru{акварелью}.
\end{quote}

\begin{quote}
\emph{Oná risúyet akvarélyu.}
\end{quote}

\begin{quote}
She paints in watercolour.
\end{quote}

Feminine.

\begin{cousinweb}
\emph{\ru{акварель}} was borrowed from French \emph{aquarelle} (EtymDB), which is usually traced to Italian \emph{acquerello}, built on \emph{acqua} 'water', from Latin \emph{aqua}.
\end{cousinweb}

\begin{grammarlens}[title={What you've built}]
One more word for telling, explaining and describing at B1.
\end{grammarlens}

\subsection*{Your turn}
\begin{itemize}
\raggedright
  \item \emph{Say it:} \emph{akvarél}
  \item \emph{Say it:} \emph{akvarél}, once more
  \item \emph{Say it:} a sentence of your own with \emph{akvarél}
  \item \emph{Recall:} say the Russian for a paintbrush, then the Russian for canvas, then say \emph{akvarél} again
\end{itemize}

\begin{tcolorbox}[breakable,colback=teal!4,colframe=teal!35!black,title={Before you move on}]
What does \ru{акварель} mean? (\textquotedblleft{}Watercolour\textquotedblright{}.) Say it once more.
\end{tcolorbox}
\section[\texorpdfstring{\ru{пейзаж} (peyzázh)}{пейзаж (peyzázh)}]{\ru{пейзаж} (peyzázh) — a landscape}

\label{lesson:RU-C521-peyzazh}



\begin{quote}
Before the new one: say the Russian for canvas, then the Russian for watercolour.
\end{quote}

\subsection*{You'll want to know: \ru{пейзаж}}
\textbf{\ru{пейзаж}} (\emph{peyzázh}) — \textquotedblleft{}a landscape\textquotedblright{}.

\begin{quote}
\ru{На} \ru{стене} \ru{висит} \ru{пейзаж} \ru{с} \ru{лесом} \ru{и} \ru{рекой}.
\end{quote}

\begin{quote}
\emph{Na stené visít peyzázh s lésom i rekóy.}
\end{quote}

\begin{quote}
A landscape with a forest and a river hangs on the wall.
\end{quote}

Masculine.

\begin{cousinweb}
\emph{\ru{пейзаж}} was borrowed from French \emph{paysage} 'countryside, landscape' (EtymDB).
\end{cousinweb}

\begin{grammarlens}[title={What you've built}]
One more word for telling, explaining and describing at B1.
\end{grammarlens}

\subsection*{Your turn}
\begin{itemize}
\raggedright
  \item \emph{Say it:} \emph{peyzázh}
  \item \emph{Say it:} \emph{peyzázh}, once more
  \item \emph{Say it:} a sentence of your own with \emph{peyzázh}
  \item \emph{Recall:} say the Russian for canvas, then the Russian for watercolour, then say \emph{peyzázh} again
\end{itemize}

\begin{tcolorbox}[breakable,colback=teal!4,colframe=teal!35!black,title={Before you move on}]
What does \ru{пейзаж} mean? (\textquotedblleft{}A landscape\textquotedblright{}.) Say it once more.
\end{tcolorbox}
\section[\texorpdfstring{\ru{позировать} (pozírovat)}{позировать (pozírovat)}]{\ru{позировать} (pozírovat) — to pose, to sit (for an artist)}

\label{lesson:RU-C521-pozirovat}



\begin{quote}
Before the new one: say the Russian for watercolour, then the Russian for a landscape.
\end{quote}

\subsection*{You'll want to know: \ru{позировать}}
\textbf{\ru{позировать}} (\emph{pozírovat}) — \textquotedblleft{}to pose, to sit (for an artist)\textquotedblright{}.

\begin{quote}
\ru{Сестра} \ru{позировала} \ru{художнику} \ru{три} \ru{часа}.
\end{quote}

\begin{quote}
\emph{Sestrá pozírovala khudózhniku tri chasá.}
\end{quote}

\begin{quote}
My sister sat for the artist for three hours.
\end{quote}

Imperfective; present \ru{позирую}. Takes the dative.

\begin{cousinweb}
\emph{\ru{позировать}} is built on \emph{\ru{поза}} 'a pose', which EtymDB traces through German \emph{Pose} to French \emph{pose}, with \emph{-\ru{ировать}}; EtymDB has no entry for the verb.
\end{cousinweb}

\begin{grammarlens}[title={What you've built}]
One more word for telling, explaining and describing at B1.
\end{grammarlens}

\subsection*{Your turn}
\begin{itemize}
\raggedright
  \item \emph{Say it:} \emph{pozírovat}
  \item \emph{Say it:} \emph{pozírovat}, once more
  \item \emph{Say it:} a sentence of your own with \emph{pozírovat}
  \item \emph{Recall:} say the Russian for watercolour, then the Russian for a landscape, then say \emph{pozírovat} again
\end{itemize}

\begin{tcolorbox}[breakable,colback=teal!4,colframe=teal!35!black,title={Before you move on}]
What does \ru{позировать} mean? (\textquotedblleft{}To pose, to sit (for an artist)\textquotedblright{}.) Say it once more.
\end{tcolorbox}
